\documentclass[11pt, a4paper]{article}
\usepackage[utf8]{inputenc}
\usepackage[margin=1in]{geometry}
\usepackage{listings}
\usepackage{xcolor}
\usepackage{tabularx}

% Configure SQL code block styling
\lstset{
    basicstyle=\ttfamily\small,
    keywordstyle=\color{blue}\bfseries,
    stringstyle=\color{red!70!black},
    commentstyle=\color{green!60!black}\itshape,
    frame=single,
    backgroundcolor=\color{gray!5},
    breaklines=true,
    showstringspaces=false
}

\title{Sheet 2 Lab Report}
\author{Mrigaj Shaw}
\date{Roll No: 2024CSB041}

\begin{document}

\maketitle

\section{Problem4}
\subsection*{Q1 : Retrieve the name of students whose name start with 'S' and contains 'r' as the second last character}
\begin{lstlisting}[language=SQL]
SELECT
  name
FROM
  students
WHERE
  TRIM(name) LIKE 'S%r_';
\end{lstlisting}
\vspace{0.3cm}
\textbf{Output:}\vspace{0.2cm}
\begin{center}
\begin{tabular}{|l|}
\hline
\textbf{name} \\ \hline
Swati Mitra                    \\ \hline
\end{tabular}
\end{center}
\newpage
\subsection*{Q2 : Retrieve the name of the youngest student(s) from the ‘CST’ department along with the total marks obtained by him (them).}
\begin{lstlisting}[language=SQL]
SELECT
  students.name,
  SUM(crs_regd.marks) as total_marks,
  AGE(CURRENT_DATE, students.bdate) as student_age
FROM
  students
  JOIN crs_regd ON crs_regd.crs_rollno = students.rollno
WHERE
  students.deptcode = 'CST'
  AND AGE(CURRENT_DATE, students.bdate) = (
    SELECT
      MIN(AGE(CURRENT_DATE, bdate))
    FROM
      students
    WHERE
      deptcode = 'CST'
  )
GROUP BY
  students.rollno,
  students.name,
  students.bdate;
\end{lstlisting}
\vspace{0.3cm}
\textbf{Output:}\vspace{0.2cm}
\begin{center}
\begin{tabular}{|l|l|l|}
\hline
\textbf{name} & \textbf{total\_marks} & \textbf{student\_age} \\ \hline
Anjali Verma                   & 176.00 & 30 years 4 mons 20 days \\ \hline
\end{tabular}
\end{center}
\newpage
\subsection*{Q3: Find age of all students}
\begin{lstlisting}[language=SQL]
SELECT
  name,
  rollno,
  AGE (CURRENT_DATE, bdate) as age
FROM
  students;
\end{lstlisting}
\vspace{0.3cm}
\textbf{Output:}\vspace{0.2cm}
\begin{center}
\begin{tabular}{|l|l|l|}
\hline
\textbf{name} & \textbf{rollno} & \textbf{age} \\ \hline
Riddhi Agarwal                 & 92005001 & 31 years 5 mons 10 days \\ \hline
Priya Verma                    & 92005002 & 30 years 1 mon 3 days \\ \hline
Rahul Das                      & 92005003 & 30 years 9 mons 15 days \\ \hline
Sneha Gupta                    & 92005004 & 30 years 7 mons 20 days \\ \hline
Ravi Kumar                     & 92005010 & 31 years 2 mons 7 days \\ \hline
Neha Patil                     & 92005011 & 29 years 10 mons 26 days \\ \hline
Vikram Singh                   & 92005012 & 30 years 8 mons \\ \hline
Ananya Roy                     & 92005013 & 30 years 4 mons 11 days \\ \hline
Suresh Nair                    & 92005100 & 31 years 17 days \\ \hline
Kavita Jain                    & 92005101 & 30 years 5 mons 28 days \\ \hline
Deepak Mehta                   & 92005102 & 31 years 3 mons 5 days \\ \hline
Pooja Reddy                    & 92005103 & 29 years 10 mons 13 days \\ \hline
Arun Bhat                      & 92005104 & 31 years 6 mons 26 days \\ \hline
Meera Iyer                     & 92005105 & 30 years 2 mons 20 days \\ \hline
Kiran Desai                    & 92005106 & 30 years 11 mons 8 days \\ \hline
Aarav Sharma                   & 92005201 & 30 years 3 mons 13 days \\ \hline
Diya Patel                     & 92005202 & 31 years 5 days \\ \hline
Kabir Mehta                    & 92005203 & 30 years 5 mons 14 days \\ \hline
Ishita Sen                     & 92005204 & 30 years 8 mons 24 days \\ \hline
Rohan Bose                     & 92005205 & 30 years 7 mons 3 days \\ \hline
Tanya Roy                      & 92005206 & 30 years 10 mons 15 days \\ \hline
Arjun Singh                    & 92005301 & 31 years 6 mons 11 days \\ \hline
Priya Das                      & 92005302 & 30 years 1 mon 16 days \\ \hline
Karan Sharma                   & 92005303 & 30 years 9 mons 5 days \\ \hline
Anjali Verma                   & 92005304 & 30 years 4 mons 20 days \\ \hline
Rahul Gupta                    & 92005305 & 30 years 11 mons 13 days \\ \hline
Sneha Patel                    & 92005306 & 30 years 7 mons \\ \hline
Vikash Kumar                   & 92005307 & 31 years 1 mon 26 days \\ \hline
Swati Mitra                    & 92005405 & 30 years 11 days \\ \hline
\end{tabular}
\end{center}
\newpage
\section{Problem5}
\subsection*{Q1 : Retrieve the name of the student(s) who obtained second highest marks in ‘DBMS’.}
\begin{lstlisting}[language=SQL]
SELECT
  s1.name
FROM
  crs_regd cr1
  JOIN students s1 ON s1.rollno = cr1.crs_rollno
  JOIN crs_offrd co ON co.crs_code = cr1.crs_cd
WHERE
  TRIM(co.crs_name) = 'DBMS'
  AND 1 = (
    SELECT
      COUNT(DISTINCT cr2.marks)
    FROM
      crs_regd cr2
      JOIN crs_offrd co2 ON co2.crs_code = cr2.crs_cd
    WHERE
      TRIM(co2.crs_name) = 'DBMS'
      AND cr2.marks > cr1.marks
  );
\end{lstlisting}
\vspace{0.3cm}
\textbf{Output:}\vspace{0.2cm}
\begin{center}
\begin{tabular}{|l|}
\hline
\textbf{name} \\ \hline
Riddhi Agarwal                 \\ \hline
\end{tabular}
\end{center}
\newpage
\subsection*{Q2: Find out the differences between the highest and lowest marks obtained in each subject}
\begin{lstlisting}[language=SQL]
SELECT
  co.crs_name,
  MAX(marks) - MIN(marks) as marks_Difference
FROM
  crs_offrd co
  JOIN crs_regd cr ON cr.crs_cd = co.crs_code
GROUP BY
  crs_name;
\end{lstlisting}
\vspace{0.3cm}
\textbf{Output:}\vspace{0.2cm}
\begin{center}
\begin{tabular}{|l|l|}
\hline
\textbf{crs\_name} & \textbf{marks\_difference} \\ \hline
Thermodynamics                      & 6.00 \\ \hline
Metallurgy                          & 0.00 \\ \hline
DBMS                                & 62.00 \\ \hline
Chemistry                           & 28.00 \\ \hline
Software Engg.                      & 0.00 \\ \hline
Aerodynamics                        & 0.00 \\ \hline
Circuit Theory                      & 12.00 \\ \hline
Physics                             & 35.00 \\ \hline
Data Structures                     & 6.00 \\ \hline
Advanced Algorithms                 & 0.00 \\ \hline
OS                                  & 16.00 \\ \hline
MIS                                 & 9.00 \\ \hline
Basic Architecture                  & 0.00 \\ \hline
Signal Processing                   & 25.00 \\ \hline
\end{tabular}
\end{center}
\newpage
\subsection*{Q3 : Assuming the existance of several interdepartmental courses, retrieve the name of the student(s) who is(are) studing under at least one faculty from each department.}
\begin{lstlisting}[language=SQL]
SELECT
  name,
  rollno
FROM
  students
  JOIN crs_regd ON crs_regd.crs_rollno = students.rollno
  JOIN crs_offrd ON crs_regd.crs_cd = crs_offrd.crs_code
  JOIN faculty on faculty.fac_code = crs_offrd.crs_fac_cd
GROUP BY
  name,
  rollno
HAVING
  COUNT(DISTINCT TRIM(faculty.fac_dept)) = (
    SELECT
      COUNT(DISTINCT TRIM(deptcode))
    FROM
      depts
  );
\end{lstlisting}
\vspace{0.3cm}
\textbf{Output:}\vspace{0.2cm}
\begin{center}
\begin{tabular}{|l|l|}
\hline
\textbf{name} & \textbf{rollno} \\ \hline
Riddhi Agarwal                 & 92005001 \\ \hline
\end{tabular}
\end{center}
\newpage
\subsection*{Q4: Assuming the existance of several interdepartmental courses, retrieve the name of the student(s) who is(are) studing under the faculties only from his(their) own department.}
\begin{lstlisting}[language=SQL]
SELECT
  name,
  rollno
FROM
  students
  JOIN crs_regd ON crs_regd.crs_rollno = students.rollno
  JOIN crs_offrd ON crs_regd.crs_cd = crs_offrd.crs_code
  JOIN faculty on TRIM(faculty.fac_code) = TRIM(crs_offrd.crs_fac_cd)
GROUP BY
  name,
  rollno,
  deptcode
HAVING
  COUNT(DISTINCT TRIM(faculty.fac_dept)) = 1
  AND MAX(TRIM(faculty.fac_dept)) = TRIM(students.deptcode)
  AND COUNT(DISTINCT TRIM(faculty.fac_code)) = (
    SELECT
      COUNT(DISTINCT TRIM(faculty.fac_code))
    FROM
      faculty
    WHERE
      TRIM(fac_dept) = TRIM(deptcode)
  );
\end{lstlisting}
\vspace{0.3cm}
\textbf{Output:}\vspace{0.2cm}
\begin{center}
\begin{tabular}{|l|l|}
\hline
\textbf{name} & \textbf{rollno} \\ \hline
Priya Verma                    & 92005002 \\ \hline
\end{tabular}
\end{center}
\newpage
\section{Problem6}
\subsection*{Q1: Display the highest parent incomes, in descending order, for each department excluding ‘ARCH’ such that only those highest parent incomes will appear that are below 12,000.}
\begin{lstlisting}[language=SQL]
SELECT
  s1.rollno,
  s1.name,
  s1.deptcode,
  s1.parent_inc
FROM
  students s1
WHERE
  deptcode != 'ARC'
  AND parent_inc < 12000
  AND s1.parent_inc = (
    SELECT
      MAX(s2.parent_inc)
    FROM
      students s2
    WHERE
      s1.deptcode = s2.deptcode
      AND s2.deptcode != 'ARC'
      AND s2.parent_inc < 12000
  )
ORDER BY
  s1.parent_inc DESC;
\end{lstlisting}
\vspace{0.3cm}
\textbf{Output:}\vspace{0.2cm}
\begin{center}
\begin{tabular}{|l|l|l|l|}
\hline
\textbf{rollno} & \textbf{name} & \textbf{deptcode} & \textbf{parent\_inc} \\ \hline
92005205 & Rohan Bose                     & IT  & 11800.0 \\ \hline
92005202 & Diya Patel                     & CSE & 11500.0 \\ \hline
92005204 & Ishita Sen                     & ELE & 9500.0 \\ \hline
\end{tabular}
\end{center}
\newpage
\subsection*{Q2: Retrieve the 5th largest parental income for hostel no. 5}
\begin{lstlisting}[language=SQL]
SELECT
  s1.rollno,
  s1.name,
  s1.parent_inc
FROM
  students s1
WHERE
  4 = (
    SELECT
      COUNT(DISTINCT s2.parent_inc)
    FROM
      students s2
    WHERE
      s2.parent_inc > s1.parent_inc
  );
\end{lstlisting}
\vspace{0.3cm}
\textbf{Output:}\vspace{0.2cm}
\begin{center}
\begin{tabular}{|l|l|l|}
\hline
\textbf{rollno} & \textbf{name} & \textbf{parent\_inc} \\ \hline
92005101 & Kavita Jain                    & 650000.0 \\ \hline
\end{tabular}
\end{center}
\newpage
\subsection*{Q3: Find the roll number of the students from each department who obtained highest total marks in their own department.}
\begin{lstlisting}[language=SQL]
SELECT
  s1.name,
  s1.rollno,
  s1.deptcode,
  SUM(cr1.marks) as total_marks
FROM
  students s1
  JOIN crs_regd cr1 ON cr1.crs_rollno = s1.rollno
GROUP BY
  s1.name,
  s1.rollno,
  s1.deptcode
HAVING
  SUM(cr1.marks) = (
    SELECT
      (MAX(dept_totals.total))
    FROM
      (
        SELECT
          s2.deptcode,
          SUM(cr2.marks) as total
        FROM
          students s2
          JOIN crs_regd cr2 ON cr2.crs_rollno = s2.rollno
        GROUP BY
          s2.rollno,
          s2.deptcode
      ) as dept_totals
    WHERE
      dept_totals.deptcode = s1.deptcode
  )
ORDER BY
  total_marks;
\end{lstlisting}
\vspace{0.3cm}
\textbf{Output:}\vspace{0.2cm}
\begin{center}
\begin{tabular}{|l|l|l|l|}
\hline
\textbf{name} & \textbf{rollno} & \textbf{deptcode} & \textbf{total\_marks} \\ \hline
Kiran Desai                    & 92005106 & MME & 71.00 \\ \hline
Ravi Kumar                     & 92005010 & ELE & 125.00 \\ \hline
Arun Bhat                      & 92005104 & IT  & 134.00 \\ \hline
Pooja Reddy                    & 92005103 & EE  & 175.00 \\ \hline
Deepak Mehta                   & 92005102 & ETC & 187.00 \\ \hline
Suresh Nair                    & 92005100 & CST & 295.00 \\ \hline
Riddhi Agarwal                 & 92005001 & CSE & 1101.00 \\ \hline
\end{tabular}
\end{center}
\newpage

\end{document}